\section{実験設定}
\label{sec:experiment}

\subsection{概要}

Unity（URP）工場シーン \texttt{FactoryExperiment} 上で，
Meta Quest 3（OpenXR）を用いたVR実験を行う．
10〜20台のAGVが荷物搬送を反復する環境で，
被験者は \textbf{Station A $\rightarrow$ Station B} へ移動する．

\begin{table}[htbp]
  \centering
  \caption{実験空間}
  \label{tab:layout}
  \begin{tabular}{ll}
    \toprule
    項目 & 値 \\
    \midrule
    Station A（開始） & $(22,\, 0.2,\, 9)$ \\
    Station B（ゴール） & $(9,\, 0.2,\, 58)$ \\
    ゴール判定 & Station B 半径 \SI{3}{m} 以内 \\
    歩行速度 & \SI{1.4}{m/s} \\
    AGV速度 & \SIrange{1.0}{2.0}{m/s}（ケースごと固定） \\
    AGV台数 & 10--20台（ケースごと固定） \\
    \bottomrule
  \end{tabular}
\end{table}

\subsection{デザイン}

\begin{itemize}[leftmargin=2em]
  \item \textbf{計画}: 被験者内，3条件 $\times$ 10ケース $=$ 30試行／人
  \item \textbf{条件}: Baseline / No-AR / Proposed
  \item \textbf{ケース固定}: 同一ケース番号では乱数シードを固定し，
        \textbf{3条件でAGVの台数・速度・軌道を完全一致}させる
  \item \textbf{被験者数}: 予備実験後に決定（パイロット $n \ge 8$ 推奨）
\end{itemize}

同一ケースでAGV挙動を揃えることは，条件間比較の\textbf{内部妥当性}を確保するための
最も重要な統制である．

\subsection{手続き}

\begin{enumerate}[leftmargin=2em]
  \item 説明・同意，HMD装着，操作練習
  \item 実験者がPCメニューで条件とケース（1--10）を選択
  \item 3秒イントロ後，Station A に出現
  \item 被験者はStation B へ移動（左スティック，HMD視点）
  \item ゴール到着で試行終了 $\rightarrow$ データ保存 $\rightarrow$ 次ケース
\end{enumerate}

\subsection{従属変数}

\begin{table}[htbp]
  \centering
  \caption{従属変数（仮説 H1--H3 との対応）}
  \label{tab:dv}
  \begin{tabular}{lllp{5.2cm}}
    \toprule
    変数 & 単位 & 検証仮説 & 定義 \\
    \midrule
    完了時間 & s & H2, H3 & A出発からB到達まで \\
    ニアミス & 回 & H1 & 閾値\SI{0.8}{m}以内接近 \\
    移動距離 & m & H2, H3 & 水平軌跡長 \\
    AGV総待機時間 & s & H1, H2 & 人が進路前方にいることによる AGV 停止時間の合計 \\
    経路上最接近距離 & m & H1 & 人が AGV 計画経路上にいるときの水平最接近 \\
    HMDヨー回転量 & ° & H2 & ケース中の絶対ヨー変化の累積 \\
    \bottomrule
  \end{tabular}
\end{table}

ニアミス閾値 \SI{0.8}{m} $=$ AGV半径\SI{0.5}{m} $+$ 被験者半径\SI{0.3}{m}．
AGVごとにクールダウン \SI{1.5}{s} を設け，同一接近の重複カウントを防ぐ．
AGVは計画経路前方\SI{1.5}{m}以内・帯半幅\SI{0.8}{m}に人がいれば目の前で停止する．
Proposed の危険度は当面 \textbf{距離のみ}（$R=R_d$）で運用し，
Inspector の \texttt{RiskScoringConfig} で Full（PTTC+CTTC+近接）へ切り替え可能とする．

\subsection{データ記録}

1セッション1 Excel（.xlsx）．
Summary に条件・ケース・従属変数，ケース別シートに \SI{0.5}{s} 間隔の位置を記録する．

\subsection{実験設定の妥当性}

本節で，\textbf{なぜこの実験設定でRQに答えられるか}を整理する．

\paragraph{構成概念の妥当性}
測りたい概念と指標の対応を明示する（\Cref{tab:construct}）．

\begin{table}[htbp]
  \centering
  \caption{構成概念の操作化}
  \label{tab:construct}
  \begin{tabular}{ll}
    \toprule
    測りたいこと & どう測るか \\
    \midrule
    移動効率 & 完了時間，移動距離，HMDヨー回転量 \\
    安全性 & ニアミス回数（\SI{0.8}{m}），経路上最接近距離 \\
    AGV側コスト & 人が原因の総待機時間 \\
    経路情報の有無 & Baseline / No-AR / Proposed \\
    危険度連動の有無 & Proposed のみ $R$ 連動（当面は距離のみ） \\
    \bottomrule
  \end{tabular}
\end{table}

H1--H3 それぞれに対応する従属変数が定義されており，
\textbf{仮説と計測が1対1で結びついている}．

\paragraph{内部妥当性（条件間比較の信頼性）}

\begin{itemize}[leftmargin=2em]
  \item \textbf{シード固定}: 同一ケースでAGV挙動を一致 $\rightarrow$ 条件差は表示方式のみ
  \item \textbf{被験者内}: 個人差（歩行速度，空間認知）を条件効果から分離
  \item \textbf{自動計測}: 観察者バイアスを低減
  \item \textbf{空間固定}: Station A/B，NavMesh，レイアウトは全条件共通
\end{itemize}

残存する脅威（条件順序の学習効果，VR酔い）には，
カウンターバランス・休憩・事前練習で対処する．

\paragraph{外部妥当性（結果の解釈範囲）}

VR工場は実工場の簡略化であるが，
\textbf{「多数AGVと作業員が通路を共有し，時間的圧力下で回避する」}
という本質的状況は再現できる．
結果は「AR経路eHMI方式の\emph{相対比較}」として解釈するのが適当である．

\paragraph{統計的分析（予定）}

\begin{itemize}[leftmargin=2em]
  \item 3条件の反復測定ANOVA（または Friedman 検定）
  \item 事後: Proposed vs.\ Baseline，Proposed vs.\ No-AR（Bonferroni補正）
  \item 効果量（$\eta^2$ または Cohen's $d$）を報告，$\alpha=0.05$
\end{itemize}

\paragraph{妥当性の総括}

\begin{table}[htbp]
  \centering
  \caption{研究全体の妥当性チェックリスト}
  \label{tab:validity}
  \begin{tabular}{p{2.8cm}p{9.8cm}}
    \toprule
    観点 & 本研究での対応 \\
    \midrule
    問いの明確さ & RQ：危険度連動表示の効果 \\
    独立変数 & 3条件で危険度エンコーディングを分離 \\
    従属変数 & 効率・安全を自動計測 \\
    統制 & シード固定でAGV一致 \\
    仮説 & H1--H3 が従属変数と対応 \\
    一般化 & VR内比較として解釈（現場検証は今後） \\
    \bottomrule
  \end{tabular}
\end{table}

以上より，本研究は
\textbf{「危険度連動eHMIが，均一表示・非表示と比べて
効率と安全のバランスを改善するか」}
という問いに，方法論的に妥当な形で答えられる計画である．
